% sections/introduction.tex

\section{Introduction}
\label{sec:introduction}

Text-to-SQL, the task of automatically translating natural language questions into executable SQL queries, has seen substantial progress with large language models (LLMs). 
In-context learning approaches such as DIN-SQL~\cite{pourreza2023din} and DAIL-SQL~\cite{gao2024dail} demonstrate strong cross-database generalization without task-specific fine-tuning, while agent-based systems like MAC-SQL~\cite{wang2025macsql} and CHESS~\cite{talaei2024chess} embed SQL generation within multi-agent pipelines that support schema-aware decomposition, iterative self-correction, and execution-guided refinement.

% 背景：Text-to-SQL 领域在 LLM 时代已取得显著进展，定义任务 +  技术路线一（ICL）： prompting 路线已经走得很远。+ 技术路线二（Agent）：进一步的执行能力


Despite this progress, a critical limitation persists: current Text-to-SQL agents cannot learn from their own execution history. 
When an agent encounters an error such as a schema misuse, an incorrect aggregation, or a malformed window function, it may fix the immediate instance through retry, but the lesson is lost once the session ends. 
The same error pattern recurs on different queries because the agent has no mechanism to retain, abstract, and reuse past experience across queries.

% 提出问题：当前 Agent 无法从自身执行历史中学习，
%   错误修复是"一次性的"，无法跨 query 复用。
% [逻辑链] P1 问题 → P2 问题现象 →  P3 问题原因。

This is not a hypothetical concern. Spider~2.0~\cite{lei2024spider2} reports that the best code agent framework achieves only 21.3\% success rate on enterprise-level Text-to-SQL workflows, revealing the extreme difficulty of real-world SQL generation. BEAVER~\cite{chen2024beaver} analyzes enterprise SQL failures and finds that analytical and structural errors (including incorrect grouping, advanced functions, and query decomposition) account for 46.8\% of failures, while schema linking errors (table selection, column selection, and join-related mistakes) contribute 35.2\%---together these two categories cover over 80\% of all errors, indicating that failures cluster in a small number of well-defined categories rather than distributing randomly.
At the same time, enterprise queries in BEAVER average 5.7 joins per query, making relation-level errors particularly prevalent and structurally repetitive. 
Taken together, these findings indicate that enterprise SQL failures are \emph{systematic}: they reflect knowledge gaps in specific schema-grounded decisions (table selection, field reference, relation paths, filtering logic) that repeat across queries until explicitly addressed.

% 核心论点：证明 P2 的 gap 不是假设性的，而是一个在企业级场景中被量化验证的严重问题；
%   并进一步揭示错误的"系统性/重复性"本质——为"可学习"提供逻辑前提。指出 gap-定量佐证 gap 的严重性 + 揭示错误的系统性本质（= 可学习前提）- 指出现有方法为何不够。

Existing memory mechanisms fall short in three specific ways. 
First, agents such as Reflexion~\cite{shinn2023reflexion} accumulate self-reflections only across retries of a \emph{single} problem instance (bounded to 1--3 entries via a sliding window); these reflections cannot transfer to novel queries, providing no cross-query knowledge accumulation~(\textbf{G1}). ExpeL~\cite{zhao2023expel} does persist experiences across tasks but stores them as free-text without SQL-specific structural constraints. 
Second, generic experience memories, whether vector embeddings or natural language summaries, lose the structural information critical for SQL generation: trigger conditions, execution steps, schema applicability, and error dimension classification~(\textbf{G2}). 
SQL's binary execution feedback (correct/incorrect) provides a natural verification signal that generic memories cannot exploit. 
Third, systems like LPE-SQL~\cite{chu2024lpesql} and Memo-SQL~\cite{yang2026memosql} maintain experience-augmented memories but lack verified consolidation: LPE-SQL logs every query outcome immediately and indiscriminately, risking noise accumulation as the knowledge base scales, while Memo-SQL pre-constructs a static experience bank offline that cannot adapt to novel error patterns at deployment~(\textbf{G3}). 
Neither implements quality-gated writing to distinguish systematic defects from one-off anomalies.

% [ 分析现有记忆机制的三个具体不足（G1/G2/G3）， 跨 query 泛化 ；结构化 ；验证性沉淀，先要"跨 query 可用"→ 再要"存得好"→ 最后要"存得对"。分别对应后文的：三层记忆架构 / 结构化 Skill 表示 / 延迟批量蒸馏。


We propose \textbf{SQLSkill}, a self-evolving Text-to-SQL agent that addresses these gaps through four core ideas (see Figure~\ref{fig:framework} for an architectural overview). 
First, we define \emph{Structured SQL Skills} as reusable workflow documents encoding trigger conditions, execution steps, applicable context, and polarity (positive or anti-pattern), going beyond code snippets~\cite{wang2023voyager} or vector embeddings. 
Second, we organize memory into \emph{three layers with evidence grading}: knowledge flows from short-term Session State through a mid-term Case Library (graded L0--L3) to a long-term Skill Bank, with promotion governed by formalized evidence conditions across four independent error dimensions (Figure~\ref{fig:evidence_flow}). 
Third, inspired by TextGrad's~\cite{yuksekgonul2024textgrad} insight that text can serve as optimizable variables, we introduce \emph{deferred batch distillation} that replaces immediate updates with batch accumulation and threshold-gated Skill writes; this design is further supported by MNL's~\cite{su2025mnl} empirical finding that batch accumulation outperforms per-instance writing on a related database QA task. 
Fourth, we introduce \emph{Anti-Pattern Skills with conflict resolution}: negative constraints coexist with positive Skills under a three-priority router that exploits LLM positional attention sensitivity~\cite{liu2024lostmiddle}.
% 提出 SQLSkill 方案，用四个核心 idea 逐一回应 P4 的三个 gap， 1）三层记忆 + 证据分级（L0-L3）；2） 延迟批量蒸馏——用批量累积 + 阈值门控替代即时写入；3 Anti-Pattern Skills + 冲突路由——负约束与正 Skill 共存，


We formalize the optimization objective as minimizing \emph{Error Recurrence Rate} (ERR), the rate at which previously seen error signatures reappear across four error dimensions, subject to maintaining baseline execution accuracy. We evaluate on Spider and BIRD benchmarks using a controlled four-condition design (No Memory $\to$ Cases Only $\to$ Train-Set Distilled Skills $\to$ Full Online Distillation) to systematically isolate each architectural layer's contribution.
% 定义评估框架——形式化优化目标（ERR）（未来也可能换为SRR）+ 消融实验设计，


In summary, this paper makes four contributions:
\begin{enumerate}[leftmargin=*,nosep]
  \item A structured SQL Skill representation operating at the query-type level, with explicit polarity markers distinguishing positive workflows from Anti-Pattern constraints (Section~\ref{sec:problem}).
  \item A three-layer memory architecture with formalized evidence grading (L0--L3) and four independent error dimensions enabling precise cross-query error clustering (Section~\ref{subsec:memory}).
  \item A deferred batch distillation mechanism with dual-track (exploratory and governance) pipelines, replacing immediate noisy writes with statistically grounded updates (Section~\ref{subsec:optimization}).
  \item A conflict resolution routing mechanism and counterfactual replay credit assignment for principled Skill lifecycle management (Sections~\ref{subsec:conflict} and~\ref{subsec:optimization}).
\end{enumerate}

% 列出四条 Contribution
% C1 → Idea 1: 结构化 SQL Skill 表示 + 极性标记（正/反模式）
%       → 回应 G2（结构化缺失）
% C2 → Idea 2: 三层记忆 + 证据分级 + 四维错误分类
%       → 回应 G1（持久化缺失）+ 部分 G2（结构化分类）
% C3 → Idea 3: 延迟批量蒸馏 + 双轨管线
%       → 回应 G3（即时写入噪声）
% C4 → Idea 4: 冲突路由 + 反事实回放信用分配
%       → 额外贡献（Skill 生命周期管理）

